% Options for packages loaded elsewhere
\PassOptionsToPackage{unicode}{hyperref}
\PassOptionsToPackage{hyphens}{url}
%
\documentclass[
]{article}
\usepackage{amsmath,amssymb}
\usepackage{iftex}
\ifPDFTeX
  \usepackage[T1]{fontenc}
  \usepackage[utf8]{inputenc}
  \usepackage{textcomp} % provide euro and other symbols
\else % if luatex or xetex
  \usepackage{unicode-math} % this also loads fontspec
  \defaultfontfeatures{Scale=MatchLowercase}
  \defaultfontfeatures[\rmfamily]{Ligatures=TeX,Scale=1}
\fi
\usepackage{lmodern}
\ifPDFTeX\else
  % xetex/luatex font selection
\fi
% Use upquote if available, for straight quotes in verbatim environments
\IfFileExists{upquote.sty}{\usepackage{upquote}}{}
\IfFileExists{microtype.sty}{% use microtype if available
  \usepackage[]{microtype}
  \UseMicrotypeSet[protrusion]{basicmath} % disable protrusion for tt fonts
}{}
\makeatletter
\@ifundefined{KOMAClassName}{% if non-KOMA class
  \IfFileExists{parskip.sty}{%
    \usepackage{parskip}
  }{% else
    \setlength{\parindent}{0pt}
    \setlength{\parskip}{6pt plus 2pt minus 1pt}}
}{% if KOMA class
  \KOMAoptions{parskip=half}}
\makeatother
\usepackage{xcolor}
\usepackage[margin=1in]{geometry}
\usepackage{longtable,booktabs,array}
\usepackage{calc} % for calculating minipage widths
% Correct order of tables after \paragraph or \subparagraph
\usepackage{etoolbox}
\makeatletter
\patchcmd\longtable{\par}{\if@noskipsec\mbox{}\fi\par}{}{}
\makeatother
% Allow footnotes in longtable head/foot
\IfFileExists{footnotehyper.sty}{\usepackage{footnotehyper}}{\usepackage{footnote}}
\makesavenoteenv{longtable}
\usepackage{graphicx}
\makeatletter
\def\maxwidth{\ifdim\Gin@nat@width>\linewidth\linewidth\else\Gin@nat@width\fi}
\def\maxheight{\ifdim\Gin@nat@height>\textheight\textheight\else\Gin@nat@height\fi}
\makeatother
% Scale images if necessary, so that they will not overflow the page
% margins by default, and it is still possible to overwrite the defaults
% using explicit options in \includegraphics[width, height, ...]{}
\setkeys{Gin}{width=\maxwidth,height=\maxheight,keepaspectratio}
% Set default figure placement to htbp
\makeatletter
\def\fps@figure{htbp}
\makeatother
\setlength{\emergencystretch}{3em} % prevent overfull lines
\providecommand{\tightlist}{%
  \setlength{\itemsep}{0pt}\setlength{\parskip}{0pt}}
\setcounter{secnumdepth}{-\maxdimen} % remove section numbering
\ifLuaTeX
  \usepackage{selnolig}  % disable illegal ligatures
\fi
\IfFileExists{bookmark.sty}{\usepackage{bookmark}}{\usepackage{hyperref}}
\IfFileExists{xurl.sty}{\usepackage{xurl}}{} % add URL line breaks if available
\urlstyle{same}
\hypersetup{
  pdftitle={COMPARISON OF DIFFERENT SQUID SIMULATIONS},
  pdfauthor={Ignacio Payá, IFOP.},
  hidelinks,
  pdfcreator={LaTeX via pandoc}}

\title{COMPARISON OF DIFFERENT SQUID SIMULATIONS}
\author{Ignacio Payá, IFOP.}
\date{mayo 26, 2024}

\begin{document}
\maketitle

{
\setcounter{tocdepth}{2}
\tableofcontents
}
\hypertarget{section}{%
\subparagraph{}\label{section}}

\hypertarget{summary}{%
\section{SUMMARY}\label{summary}}

\begin{center}\rule{0.5\linewidth}{0.5pt}\end{center}

This report was developed in rmakdown to compare results of different
cases of Humboldt Squid Stock simulations. It collects the output files
of ``Run HSquid\_Rmd.R'' script. The parameters by cases are shown in
tables and the main variables (recruitments, biomass, fishing
selectivity, catches, etc.) are shown in graphs.

\hypertarget{introduction}{%
\section{INTRODUCTION}\label{introduction}}

\begin{center}\rule{0.5\linewidth}{0.5pt}\end{center}

The `HSquid\_Html.rmd' rmardwon script was developed to simulate
Humboldt Squid Stock dynamics using an input parameter file named
HSquid\_Par.csv.

The `Run HSquid\_Rmd.R' R script was written to conduct different
simulation by means of run `HSquid\_Html.rmd' with different input
parameters and to save parameters and main results in csv format file
and graphs in png format files.

Therefore, it was neccesary to write a rmarkdown script to report on the
different simulations outputs (`Run Squid Rmd.R').

\hypertarget{methods}{%
\section{METHODS}\label{methods}}

\begin{center}\rule{0.5\linewidth}{0.5pt}\end{center}

No calculations were done in this script. The figures and tables are
read from the current folder.

The individual, population and fishery parameters are shown in tables 1
to 3.

\hypertarget{section-1}{%
\subparagraph{}\label{section-1}}

\hypertarget{results}{%
\section{RESULTS}\label{results}}

\begin{center}\rule{0.5\linewidth}{0.5pt}\end{center}

\hypertarget{table-1.-individual-parameters-by-case.}{%
\subsection{Table 1. Individual Parameters by
case.}\label{table-1.-individual-parameters-by-case.}}

\begin{longtable}[]{@{}rrrrrrr@{}}
\toprule\noalign{}
Case & a\_expo & b\_expo & a & b & lm50 & lmrange \\
\midrule\noalign{}
\endhead
\bottomrule\noalign{}
\endlastfoot
1 & 309.11 & 0.0029 & 2.31e-05 & 3.077 & 80 & 10 \\
2 & 309.11 & 0.0029 & 2.31e-05 & 3.077 & 80 & 10 \\
3 & 309.11 & 0.0029 & 2.31e-05 & 3.077 & 80 & 10 \\
\end{longtable}

\hypertarget{table-2.-season-recruitment-parameters-by-case.}{%
\subsection{Table 2. Season Recruitment Parameters by
case.}\label{table-2.-season-recruitment-parameters-by-case.}}

\begin{longtable}[]{@{}rrr@{}}
\toprule\noalign{}
Case & RSeason & SdSeason \\
\midrule\noalign{}
\endhead
\bottomrule\noalign{}
\endlastfoot
1 & 1 & 0 \\
2 & 3 & 0 \\
3 & 3 & 0 \\
\end{longtable}

\hypertarget{table-3.-population-parameters-by-case.}{%
\subsection{Table 3. Population Parameters by
case.}\label{table-3.-population-parameters-by-case.}}

\begin{longtable}[]{@{}rrrrrrrrr@{}}
\toprule\noalign{}
Case & M\_Annual & M\_bin & M\_cv & SB0 & SRModel & h & prev.rsigma &
rsigma \\
\midrule\noalign{}
\endhead
\bottomrule\noalign{}
\endlastfoot
1 & 1.5 & 0.125 & 0 & 1e+07 & 1 & 0.9 & 0.0 & 0.0 \\
2 & 1.5 & 0.125 & 0 & 1e+07 & 1 & 0.9 & 0.0 & 0.0 \\
3 & 1.5 & 0.125 & 0 & 1e+07 & 1 & 0.9 & 0.1 & 0.1 \\
\end{longtable}

\hypertarget{table-4.-fishery-parameters-by-case.}{%
\subsection{Table 4. Fishery Parameters by
case.}\label{table-4.-fishery-parameters-by-case.}}

\begin{longtable}[]{@{}rrrrrrrrr@{}}
\toprule\noalign{}
Case & Fref1 & as1 & bs1 & cs1 & Fref2 & as2 & bs2 & cs2 \\
\midrule\noalign{}
\endhead
\bottomrule\noalign{}
\endlastfoot
1 & 0.0 & 0.00 & 0.0 & 0.00 & 0.0 & 0.00 & 0.0 & 0.00 \\
2 & 0.1 & 0.15 & 0.4 & 0.15 & 0.1 & 0.15 & 0.8 & 0.15 \\
3 & 0.2 & 0.15 & 0.5 & 0.15 & 0.1 & 0.15 & 0.9 & 0.15 \\
\end{longtable}

\hypertarget{figure-1.-n-at-first-month-by-case.}{%
\subsection{Figure 1. N at first Month by
case.}\label{figure-1.-n-at-first-month-by-case.}}

\includegraphics{CasesFig_N1.png}

\hypertarget{figure-2.-recruitment-by-case.}{%
\subsection{Figure 2. Recruitment by
case.}\label{figure-2.-recruitment-by-case.}}

\includegraphics{CasesFig_R.png}

\hypertarget{figure-3.-whole-biomass-by-case.}{%
\subsection{Figure 3. Whole Biomass by
case.}\label{figure-3.-whole-biomass-by-case.}}

\includegraphics{CasesFig_BT.png}

\hypertarget{figure-4.-spawning-biomass-by-case.}{%
\subsection{Figure 4. Spawning Biomass by
case.}\label{figure-4.-spawning-biomass-by-case.}}

\includegraphics{CasesFig_SBT.png}

\hypertarget{figure-5.-stock-versus-recruitment-by-case.}{%
\subsection{Figure 5. Stock versus Recruitment by
case.}\label{figure-5.-stock-versus-recruitment-by-case.}}

\includegraphics{CasesFig_SB_R.png}

\hypertarget{figure-6.-fleet-1-selectivity-by-case.}{%
\subsection{Figure 6. Fleet 1 Selectivity by
case.}\label{figure-6.-fleet-1-selectivity-by-case.}}

\includegraphics{CasesFig_S1.png}

\hypertarget{figure-7.-fleet-2-selectivity-by-case.}{%
\subsection{Figure 7. Fleet 2 Selectivity by
case.}\label{figure-7.-fleet-2-selectivity-by-case.}}

\includegraphics{CasesFig_S2.png}

\hypertarget{figure-8.-fleet-1-catch-tons-by-case.}{%
\subsection{Figure 8. Fleet 1 Catch (tons) by
case.}\label{figure-8.-fleet-1-catch-tons-by-case.}}

\includegraphics{CasesFig_YT1.png}

\hypertarget{figure-9.-fleet-2-catch-tons-by-case.}{%
\subsection{Figure 9. Fleet 2 Catch (tons) by
case.}\label{figure-9.-fleet-2-catch-tons-by-case.}}

\includegraphics{CasesFig_YT2.png}

\hypertarget{figure-8.-whole-catch-tons-by-case.}{%
\subsection{Figure 8. Whole Catch (tons) by
case.}\label{figure-8.-whole-catch-tons-by-case.}}

\includegraphics{CasesFig_YT.png}

\end{document}
